% TurboQuant at the Edge -- IEEE AIoT 2026 (Track 3: Edge/Cloud/Fog -- LLM deployment)
% IEEE conference template (IEEEtran). Compile with pdflatex (IEEEtran.cls from texlive/Overleaf).
%
% STATUS: skeleton with REAL established numbers + the analytical memory budget
% from benchmarks/benchmark_edge.py. Cells marked \tbd{} MUST be filled from
% on-device runs (Jetson Orin, RPi 5, a consumer GPU) before submission:
%   python benchmarks/benchmark_edge.py --models llama-3.2-1b llama-3.2-3b qwen2.5-7b \
%       --contexts 2048 8192 32768 --bits 3 --gen 128 --energy --json out/edge_<device>.json
\documentclass[conference]{IEEEtran}
\IEEEoverridecommandlockouts

\usepackage{times}
\usepackage{amsmath,amssymb}
\usepackage{graphicx}
\usepackage{listings}
\usepackage{hyperref}
\usepackage{url}
\usepackage{cite}
\usepackage{booktabs}
\usepackage{xcolor}

% Placeholder macro for numbers that must come from on-device runs.
\newcommand{\tbd}[1][TBD]{\textcolor{red}{\textbf{[#1]}}}

\begin{document}

\title{TurboQuant at the Edge: PCA--Matryoshka Quantization for\\
Foundation-Model Inference on Resource-Constrained AIoT Devices}

\author{
\IEEEauthorblockN{Andrew H. Bond, Senior Member, IEEE}
\IEEEauthorblockA{Department of Computer Engineering\\
San Jose State University, San Jose, CA, USA\\
andrew.bond@sjsu.edu}
}

\maketitle

\begin{abstract}
Deploying large language and embedding models on AIoT edge hardware is
constrained less by compute than by \emph{memory and energy}: model weights and
the autoregressive key--value (KV) cache exceed the budgets of consumer GPUs,
edge accelerators, and CPUs. We present \textbf{TurboQuant}, a unified
compression pipeline---PCA--Matryoshka dimension reduction, a random orthogonal
rotation, Lloyd--Max scalar quantization with learned codebooks, and
bit-packing---applied across the three surfaces that dominate edge memory:
\emph{model weights}, the \emph{RoPE-aware LLM KV cache}, and the \emph{embedding
store} used for on-device retrieval. On production embeddings TurboQuant attains
up to $27\times$ compression at $99.8\%$ recall@10 (with oversampling and
reranking), and learned codebooks reduce quantization error by $22\%$; at
3-bit, KV-cache memory shrinks $5.3\times$ versus fp16. We show that this lifts a
7B-parameter model from \emph{not fitting} an 8\,GB device under fp16 to fitting
comfortably under TurboQuant. We evaluate end-to-end deployment on Jetson Orin,
Raspberry Pi~5, and a consumer GPU, reporting peak memory, throughput, and
\emph{energy-per-token}. We emphasize task-relevant metrics over proxy measures:
our data shows cosine similarity is \emph{not} a reliable predictor of retrieval
quality under aggressive compression. TurboQuant is open-source (PyPI; Zenodo
DOI~10.5281/zenodo.20660087) with 397 tests, and runs on Volta-class GPUs and CPU.
\end{abstract}

\begin{IEEEkeywords}
edge AI, AIoT, LLM deployment, model compression, KV cache, quantization,
energy efficiency, on-device inference
\end{IEEEkeywords}

\section{Introduction}
Foundation models are moving to the edge---for privacy, latency, and offline
operation---but on IoT-class hardware the binding constraint is \emph{memory and
energy}, not arithmetic. Three sinks dominate: the model \emph{weights}; the
\emph{KV cache}, which grows with context length $\times$ layers and quickly
dwarfs the weights at long context; and, for retrieval-augmented edge agents,
the \emph{embedding/index store}. Prior compression work typically optimizes one
surface in isolation and reports proxy metrics.

This paper makes four contributions: (1) a \emph{single} quantization pipeline
spanning weights, KV cache, and embeddings; (2) a RoPE-aware KV-cache compressor
for long-context on-device inference; (3) an \texttt{AutoConfig}-driven
deployment path for Volta+/CPU targets; and (4) an edge evaluation reporting
memory, throughput, and energy-per-token, with the explicit caveat that cosine
similarity is not a reliable proxy for task quality.

\section{Background and Related Work}
\textbf{Weight quantization.} GPTQ~\cite{gptq} and AWQ~\cite{awq} push LLM weights
to 3--4 bits with small accuracy loss. \textbf{KV-cache quantization.}
KVQuant~\cite{kvquant} and KIVI~\cite{kivi} compress the cache that dominates
long-context memory. \textbf{Matryoshka representations}~\cite{mrl} enable
truncatable embeddings. \textbf{Vector quantization} underpins approximate
nearest-neighbor search~\cite{faiss}. The gap: edge deployment needs a
\emph{whole-system} memory budget across all three surfaces, evaluated on
task-grounded rather than proxy metrics.

\section{The TurboQuant Pipeline}
TurboQuant applies one pipeline to each memory surface
(Fig.~\ref{fig:pipeline}): (i) PCA--Matryoshka rotate-and-truncate; (ii) a random
orthogonal (QR/structured) rotation that flattens the coordinate-wise dynamic
range; (iii) a Lloyd--Max scalar quantizer with \emph{learned} codebooks
($-22\%$ quantization error vs.\ uniform); and (iv) bit-packing (e.g.,
$8\times3$-bit $\rightarrow$ 3 bytes) with the L2 norm preserved alongside the
packed bits. For the KV cache, a RoPE-aware quantizer and a hot-window of recent
tokens in full precision preserve attention fidelity; for retrieval, a
compressed HNSW index supports on-device search. Scalar quantization plus a
cheap rotation (no large codebook tables) is well suited to CPU and edge
accelerators.

\begin{figure}[t]
\centering
% \includegraphics[width=\columnwidth]{figs/pipeline.pdf}  % TODO: export from poster-assets
\fbox{\parbox{0.9\columnwidth}{\centering\vspace{1.2em} Raw $\rightarrow$
PCA--Matryoshka $\rightarrow$ orthogonal rotation $\rightarrow$ Lloyd--Max
quantize $\rightarrow$ bit-pack \vspace{1.2em}}}
\caption{The TurboQuant pipeline, applied to weights, KV cache, and embeddings.}
\label{fig:pipeline}
\end{figure}

\section{Edge Deployment Design}
We model the device memory budget as
$M_{\text{dev}} \ge M_{\text{weights}} + M_{\text{KV}}(L, B) + M_{\text{index}}$,
where $L$ is context length and $B$ batch. TurboQuant reduces each term:
weights via low-bit packing, KV via the RoPE-aware compressor (linear in layers
and context), and the index via PCA-truncation plus scalar quantization.
\texttt{AutoConfig.from\_pretrained} selects a bit-width per device tier (Volta+
GPU $\rightarrow$ edge accelerator $\rightarrow$ CPU-only).

\section{Experimental Setup}
\textbf{Devices.} \tbd[Jetson Orin Nano 8\,GB], \tbd[Jetson Orin NX 16\,GB],
\tbd[Raspberry Pi~5 8\,GB (CPU)], and a Volta-class consumer GPU.
\textbf{Models.} Llama-3.2-1B/3B, Qwen2.5-7B, Mistral-7B.
\textbf{Tasks.} a retrieval set (recall@10) and a generation task
(task accuracy / perplexity). \textbf{Baselines.} fp16, GPTQ, AWQ, and KV-cache
quantization. \textbf{Metrics.} peak memory, throughput (tok/s), and
\emph{energy-per-token} (J/tok), the last measured via on-board power telemetry
(NVML on GPU; \texttt{tegrastats} on Jetson). Memory budgets below are analytical
(2\,B/elem fp16 vs.\ bit-packed at the target width) and verified against the
library's measured footprint; throughput and energy are measured on device.

\section{Results}

\subsection{Memory budget and device fit}
Table~\ref{tab:budget} reports the weights-plus-KV budget at 8K context. At
3-bit, KV shrinks $5.3\times$ and total footprint $\approx 5\times$. The
practical consequence (Table~\ref{tab:fit}): a 7B model does \emph{not} fit an
8\,GB edge device under fp16 but fits under TurboQuant, as do all tested models.

\begin{table}[t]
\caption{Memory footprint (GB), weights $+$ KV @ 8K context (analytical).}
\label{tab:budget}
\centering
\begin{tabular}{lrrrrr}
\toprule
Model & Params & W fp16 & W 3-bit & KV fp16 & Total 3-bit \\
\midrule
Llama-3.2-1B & 1.24B & 2.31 & 0.43 & 0.25 & 0.48 \\
Llama-3.2-3B & 3.21B & 5.98 & 1.12 & 0.88 & 1.28 \\
Qwen2.5-7B   & 7.6B  & 14.16 & 2.66 & 0.44 & 2.74 \\
\bottomrule
\end{tabular}
\end{table}

\begin{table}[t]
\caption{Models that fit a device budget (TurboQuant 3-bit vs.\ fp16),
weights $+$ KV @ 8K context.}
\label{tab:fit}
\centering
\begin{tabular}{lrcc}
\toprule
Device tier & Budget & fp16 & TurboQuant \\
\midrule
Jetson Orin Nano & 8\,GB  & 2/3 & \textbf{3/3} \\
Jetson Orin NX   & 16\,GB & 3/3 & 3/3 \\
Raspberry Pi 5   & 8\,GB  & 2/3 & \textbf{3/3} \\
Consumer GPU     & 12\,GB & 2/3 & \textbf{3/3} \\
\bottomrule
\end{tabular}
\end{table}

\subsection{On-device throughput and energy}
Table~\ref{tab:device} reports measured decode throughput and energy-per-token
on each target. \emph{(To be filled from \texttt{benchmark\_edge.py --energy}
runs on hardware; numbers below are placeholders.)}

\begin{table}[t]
\caption{Measured decode throughput and energy-per-token (TurboQuant 3-bit).}
\label{tab:device}
\centering
\begin{tabular}{llrr}
\toprule
Device & Model & tok/s & J/tok \\
\midrule
Jetson Orin Nano & Qwen2.5-7B & \tbd & \tbd \\
Jetson Orin NX   & Qwen2.5-7B & \tbd & \tbd \\
Raspberry Pi 5   & Llama-3.2-3B & \tbd & \tbd \\
Consumer GPU     & Qwen2.5-7B & \tbd & \tbd \\
\bottomrule
\end{tabular}
\end{table}

\subsection{Embedding compression and retrieval quality}
For the on-device retrieval surface, TurboQuant achieves up to $27\times$
embedding compression at $99.8\%$ recall@10 with $5\times$ oversampling and
reranking, benchmarked identically across methods on 50K production
embeddings~\cite{turboquant}. Crucially, cosine similarity to the original vector
is \emph{not} a reliable proxy: PCA-256+TQ3 has lower cosine (0.963) yet higher
recall@10 (78.2\%) than PCA-384+TQ3 (0.979 cosine, 76.4\% recall). We therefore
report task-relevant retrieval metrics throughout.

\section{Discussion and Limitations}
The measured KV footprint carries fixed overhead (the full-precision hot window,
codebooks) that amortizes only at long context; for very short contexts,
uncompressed caching can be smaller. The smallest devices still cannot host the
largest models even compressed. Reranking adds compute at the edge that must be
budgeted against its retrieval-quality gain. Finally, the throughput reported by
the KV-path microbenchmark is not end-to-end model latency; full-model numbers
(Table~\ref{tab:device}) come from device runs with the model runtime.

\section{Conclusion}
TurboQuant compresses the three memory surfaces that block foundation-model
inference at the edge, turning models that overflow IoT-class budgets into ones
that fit---with task-grounded quality and on-device energy accounting. Code,
tests, and benchmarks (including \texttt{benchmark\_edge.py}) are open-source.

\section*{Reproducibility}
TurboQuant: PyPI \texttt{turboquant-pro}; DOI~10.5281/zenodo.20660087. The edge
budget/throughput/energy numbers are produced by
\texttt{benchmarks/benchmark\_edge.py}.

\begin{thebibliography}{00}
\bibitem{gptq} E. Frantar, S. Ashkboos, T. Hoefler, and D. Alistarh, ``GPTQ:
Accurate post-training quantization for generative pre-trained transformers,''
in \emph{ICLR}, 2023.
\bibitem{awq} J. Lin \emph{et al.}, ``AWQ: Activation-aware weight quantization
for LLM compression and acceleration,'' in \emph{MLSys}, 2024.
\bibitem{kvquant} C. Hooper \emph{et al.}, ``KVQuant: Towards 10 million context
length LLM inference with KV cache quantization,'' in \emph{NeurIPS}, 2024.
\bibitem{kivi} Z. Liu \emph{et al.}, ``KIVI: A tuning-free asymmetric 2-bit
quantization for KV cache,'' in \emph{ICML}, 2024.
\bibitem{mrl} A. Kusupati \emph{et al.}, ``Matryoshka representation learning,''
in \emph{NeurIPS}, 2022.
\bibitem{faiss} J. Johnson, M. Douze, and H. J\'egou, ``Billion-scale similarity
search with GPUs,'' \emph{IEEE Trans. Big Data}, 2021.
\bibitem{turboquant} A. H. Bond, ``TurboQuant Pro,'' Zenodo, 2026.
doi:10.5281/zenodo.20660087.
\end{thebibliography}

\end{document}
